
\def\PUDcodigo{ECA218}

\def\PUDcodacad{04507.56}

\def\PUDdisciplina{Identificação de Sistemas}

\def\PUDsemestre{Opt}

\def\PUDhoras{80}

%NOVOS ---------------------------------------------------
\def\PUDhorasTeoria{64}
\def\PUDhorasPratica{16}

\def\PUDinterECA{ 
\hyperref[PUD:ECA233]{Controle I (04507.31)}

\hyperref[PUD:ECA202]{Controle II (04507.40)}

\hyperref[PUD:EME132]{Inteligência Computacional Aplicada (04507.61)}
}

\def\PUDatuaisECA{- Métodos de Identificação Robusta e \textit{Online}}

\def\PUDrecursos{Projetor multimídia; Quadro branco e pincel.}

%----------------------------------------------------------

\def\PUDcreditos{4}

\def\PUDprerequisito{ \hyperref[PUD:ECA233]{ECA233} }

\def\PUDpreacad{ \hyperref[PUD:ECA233]{Controle I (04507.31)} }

\def\PUDobjetivos{Conhecer e utilizar as principais técnicas lineares e não lineares relacionadas à modelagem, estimação e identificação de sistemas dinâmicos.
%Aprender os conceitos básicos relacionados a modelagem, estimação e identificação de sistemas. Compreender as principais técnicas lineares e não-lineares de identificação de sistemas. Aplicar técnicas de estimação de parâmetros e identificação em sistemas reais.
}

\def\PUDementa{
Introdução a Identificação de Sistemas; Métodos Determinísticos e Não-Paramétricos; Estimador de Mínimos Quadrados; Propriedades Estatísticas de Estimadores; Estimadores Não Polarizados; Seleção de Estruturas e Validação de Modelos; Identificação de Ssitemas Não Lineares. 
%Modelagem Matemática;  Representações Lineares;  Métodos Determinísticos; Métodos Não-Paramétricos;  Estimador de Mínimos Quadrados; Propriedades Estatísticas de Estimadores;  Estimadores Não-Polarizados; Estimadores Recursivos.
}

\def\PUDconteudo{ & Introdução a Identificação de Sistemas: modelagem matemática, o problema de identificação de sistemas, conceitos básicos em modelagem, passo a passo em um problema de identificação de sistemas, noções de processos aleatórios. 
%Modelagem Matemática: Introdução.  Conceitos Básicos.  Estimação de Parâmetros.  Identificação de Sistemas.  Simulação de Modelos. 
\\ 
 & Métodos Determinísticos e Não Paramétricos: alguns casos simples, método de Sundaresan, identificação em malha fechada, identificação usando convolução, reduzindo o efeito do ruído no domínio do tempo, sinais aleatórios e pseudo-aleatórios.
 %Representações Lineares: Introdução.  Funções de Transferência.  Resposta Temporal.  Resposta em Frequência.  Representação no Espaço de Estados.  Representações em Tempo Discreto. 
\\ 
 & Estimador de Mínimos Quadrados: sistemas de equações lineares, algoritmo de mínimos quadrados ordinários (OLS), filtragem adaptativa linear, filtro de Wiener, algoritmo Least Mean Squares (LMS), algoritmo Recursive Least Squares (RLS), ortogonalidade do estimador OLS.
 %Métodos Determinísticos: Introdução.  Alguns Casos Simples.  O Método de Sundaresan.  Identificação em Malha Fechada.  Identificação Usando Convolução.  Identificação no Domínio da Frequência. 
\\ 
 & Propriedades Estatísticas de Estimadores: introdução, conceito de polarização, polarização do estimador OLS, interpretação de polarização, polarização em modelos ARX, o problema de erros nas variáveis, covariância de estimadores, variância mínima do estimador OLS.
 %Métodos Não-Paramétricos: Introdução.  Reduzindo o Efeito do Ruído no Domínio do Tempo.  Sinais Aleatórios e Pseudo-Aleatórios.  Reduzindo o Efeito do Ruído no Domínio a Frequência.  Persistência de Excitação.  
\\ 
 & Estimadores Não Polarizados: introdução, estimador Extended Least Squares (ELS), estimador Generalized Least Squares (GLS), estimador Instrumental Variables (IV).
 %Estimador de Mínimos Quadrados: Introdução.  Gerando o Sistema de Equações.  O Método de Mínimos Quadrados.  Propriedades.  Estimação de Parâmetros de Modelos ARX Usando Mínimos Quadrados. 
\\ 
 & Seleção de Estruturas e Validação de Modelos: introdução, seleção da ordem de modelos lineares, análise de resíduos, outras métricas de validação.
 %Propriedades Estatísticas de Estimadores: Introdução.  Polarização de Estimadores.  Polarização do Estimador de Mínimos Quadrados. Covariância de Estimadores.  Eficiência de Estimadores. 
\\ 
 & Identificação de Sistemas Não Lineares: algumas representações NARMAX, estimador de mínimos quadrados ortogonais, método clássico de Gram-Schmidt (CGS), método de Golub-Householder (GH), estimadores baseados em inteligência artificial. 
 %Estimadores Não Polarizados: Introdução.  Estimador de Mínimos Quadrados Estendido.  Estimador de Mínimos Quadrados Generalizado. Método das Variáveis Instrumentais. 
\\ 
 & Aplicações: exemplos de aplicações práticas usando modelos lineares baseados em mínimos quadrados, modelos não lineares baseados em mínimos quadrados, modelos não lineares usando redes neurais artificias e métodos de kernel.
 %Estimadores Recursivos: Introdução.  Atualização Recursiva.  Estimador de Mínimos Quadrados Recursivo.  Outros Estimadores Recursivos. Estimação de Parâmetros Variantes no Tempo. 
 }

\def\PUDmetodo{Aulas teóricas expositivas com auxílio de recursos audio-visuais e aulas práticas com simulações computacionais.
%Aulas expositórias que podem ser teóricas e/ou práticas, onde as práticas no laboratório serão marcadas no decorrer da disciplina.
}

\def\PUDavalia{As avaliações de conhecimento poderão ser através de provas escritas teóricas, como também através de trabalhos/projetos a serem apresentados na forma de relatórios e/ou seminários.
%A avaliação será feita com aplicação de provas teóricas e/ou práticas, além da possibilidade de inclusão de trabalhos e seminários no decorrer da disciplina.
}

\def\PUDbibbasica{ &  \href{www.biblioteca.ifce.edu.br/index.asp?codigo_sophia=65213}{Aguirre}, Luis Antonio.  Introdução à Identificação de Sistemas - Técnicas Lineares e Não-Lineares Aplicadas a Sistemas Reais.  4ª ed.  Belo Horizonte: Editora UFMG.  2015. 774p.  ISBN: 9788542300796.
\\ 
 &  \href{www.biblioteca.ifce.edu.br/index.asp?codigo_sophia=64121}{Billings}, Stephen A.  Nonlinear System Identification: NARMAX Methods in the Time, Frequency and Spatio-Temporal Domains.  1 ed.  Wiley.  2013.  555p.  ISBN: 9781119943594. 
\\ 
 &  \href{www.biblioteca.ifce.edu.br/index.asp?codigo_sophia=63228}{Ogata}, Katsuhiko.  Engenharia de controle moderno.  5º ed.  São Paulo, SP: Pearson Prentice Hall, 2010.  809p.  ISBN: 9788576058106. }
%  LJUNG, Lennart;  System Identification: Theory for the User.  2 ed.  New Jersey: Prentice Hall.  1999.  609p.  ISBN: 9780136566953.

\def\PUDbibcomplementar{ &  \href{www.biblioteca.ifce.edu.br/index.asp?codigo_sophia=62576}{Maya}, Paulo Álvaro.  Controle essencial.  2º ed.  São Paulo, SP: Pearson Education do Brasil, 2014.  347p.  ISBN: 9788543002415.
%  ISERMAN, R.;  MUNCHHOF, M.;  Identification of Dynamic Systems - An Introduction With Applications.  1 ed.   London: Springer.  2010.  705p.  ISBN: 9783540788782. 
\\ 
 & \href{biblioteca.ifce.edu.br/index.asp?codigo_sophia=63123}{Nise}, Norman S. Engenharia de sistemas de Controle. 6º ed. Rio de Janeiro, RJ: LTC, 2012. 745p. ISBN: 9788521621355.
 %\href{www.biblioteca.ifce.edu.br/index.asp?codigo_sophia=24551}{Almeida}, José Luiz Antunes de.  Dispositivos semicondutores: tiristores : controle de potência em CC e CA.  12º ed.  São Paulo, SP: Érica, 2013.  150p.  ISBN: 9788571942981.
%  SODERSTROM, T.;  STOICA, P.;  System Identification.  1 ed.  Prentice Hall International.  2001.  646p.  ISBN: 9780138812362. }
\\
 & \href{biblioteca.ifce.edu.br/index.asp?codigo_sophia=24230}{Campos}, Mario Cesar M. Massa. Controles típicos de equipamentos e processos industriais. Volume único. Editora Edgar Blucher. São Paulo, 2006. ISBN: 8521203985.
 %\href{www.biblioteca.ifce.edu.br/index.asp?codigo_sophia=41951}{Macintyre}, Archibald Joseph.  Instalações hidráulicas: prediais e industriais.  4º ed.  Rio de Janeiro, RJ: LTC, 2010.  579p.  ISBN: 9788521616573.
\\
 & \href{biblioteca.ifce.edu.br/index.asp?codigo_sophia=65382}{Capelli}, Alexandre. Automação industrial: controle do movimento e processos contínuos. 3º ed. São Paulo, SP: Érica, 2013. 236p. ISBN: 9788536501178.
 %\href{www.biblioteca.ifce.edu.br/index.asp?codigo_sophia=55455}{Rosário}, João Maurício. E-book. Pearson. 362p. ISBN: 9788576050100. 
 \\
 & \href{www.biblioteca.ifce.edu.br/index.asp?codigo_sophia=24226}{Simões}, Marcelo Godoy.  Controle e modelagem fuzzy.  2º ed.  São Paulo, SP: Blucher: FAPESP, 2007.  ISBN: 9788521204169.
%  COELHO, Antonio A.  R.;  COELHO, Leandro S.;  Identificação de Sistemas Dinâmicos Lineares.  1 ed.  Florianópolis: Editora UFSC.  2004.  181p.  ISBN: 9788532802804.
}

\def\PUDautor{José Daniel de Alencar Santos}

\def\PUDdata{2013-05-19}

\def\PUDrevisao{2}

\def\PUDrevisor{José Daniel de Alencar Santos}

\def\PUDdatarevisao{2019-05-15}

\PUDcorpo
